\documentclass[10pt]{article}
\usepackage{icmlko}

\icmlmeta{IP 파운데이션 월드모델}{73}{게이트를 먹인 것은 마흔둘이 아니라 여섯이었다}{출처들끼리 라벨 없이 서로의 방향을 알아보지만, 도메인을 건너 가중치를 배우는 것은 아직 이르다}

\begin{document}

\twocolumn[
\icmlheader

\begin{abstract}
방향에 라벨이 든다는 것을 노트 70이 정했고(32건) 노트 72가 그 이유를 밝혔다.
그런데 안 쓴 정보가 하나 있었다 --- \textbf{다른 출처들의 의견}. 팝업을 맞히는
여섯 중 다섯이 $+$0.4대로 같은 쪽이고 애니만 $-$0.31이면, 애니가 뒤집혔다는
것을 \textbf{팝업 라벨을 한 건도 안 보고} 알 수 있다. 부호를 둘로 쪼갠다 ---
\textbf{상대 방향}(출처들끼리 맞추기, 라벨 0건)과 \textbf{전역 방향}(맞춰진
전체가 어느 쪽인지, 라벨 필요). 합의 정향은 \textbf{일곱 대상 전부에서 애니를
정확히 지목한다}(팝업 가중치: 다섯이 $+$0.83$\sim+$1.00, 애니 $-$0.799). 그런데
\textbf{고쳐도 판정치가 안 는다} --- $+$0.0073 {[}$-$0.0035, $+$0.0214{]}, 보류.
노트 72가 밝힌 희석 때문이고, 대상별로 부호가 갈린다(아이돌 $+$0.064 이득,
\textbf{팝업 $-$0.024 손해}). 팝업에서는 뒤집힌 출처를 \emph{그냥 두는 것}이
뒤집거나 빼는 것보다 낫다 --- 음의 상관 구성원이 상관된 오차를 상쇄한다. 얻은
것은 성적이 아니라 \textbf{단순함}이다: 결정 여섯을 하나로 줄여도 같은 자리다
($k{=}32$에서 $+$0.3327 대 $+$0.3367). 그리고 \textbf{게이트}를 만들었다 ---
쌍의 성질로 셀 성적을 예측해 가중치를 정하는, 대상 라벨 0건짜리 혼합 전문가.
\textbf{진다}($+$0.2473 대 균등 $+$0.2744). 진단이 정확하다 --- \textbf{대상
수준 특징은 대상 안에서 상수라 유효 학습 표본이 셀 42개가 아니라 도메인 6개다.}
빼면 게이트가 합의로 되돌아간다($+$0.2836). 여덟째 도메인의 값이 여기서
분명해진다 --- 셀 42$\to$56, 게이트 유효 표본 6$\to$7.
\end{abstract}
]

\begin{figure*}[t]
\centering
\includegraphics[width=\textwidth]{figs/gate.pdf}
\caption{왼쪽 --- 방향 결정의 값. 오른쪽 --- 게이트의 특징 집합별 판정치.}
\end{figure*}

\section{부호를 둘로 쪼갠다}

\begin{quote}\fontsize{9.5}{11.5}\selectfont
\textbf{상대 방향} 출처들끼리 같은 쪽을 보게 맞추는 것 --- \textbf{라벨 0건}\\
\textbf{전역 방향} 맞춰진 전체가 어느 쪽인지 --- 라벨이 든다
\end{quote}

지금까지 둘을 구분하지 않아 여섯 결정을 전부 라벨로 샀다. 상대 방향은 각 출처를
뺀 나머지의 순위 평균과 비교해 반대면 뒤집는 것으로, 수렴할 때까지 반복한다.
전역 부호는 다수결로 고정한다 --- 대부분의 출처가 옳게 정향돼 있다는 가정이다.

\begin{table}[t]
\centering\tsize
\caption{팝업을 맞히는 여섯의 합의 가중치. 팝업 라벨 0건.}
\begin{tabular}{lrr}
\toprule
출처 & 합의 가중치 & 실제 $\rho$ \\
\midrule
게임 & $+$0.997 & $+$0.445 \\
도서 & $+$0.995 & $+$0.450 \\
펀딩 & $+$0.883 & $+$0.473 \\
아이돌 & $+$0.863 & $+$0.476 \\
웹툰 & $+$0.826 & $+$0.324 \\
\textbf{애니} & \textbf{$-$0.799} & \textbf{$-$0.310} \\
\bottomrule
\end{tabular}
\end{table}

\claimx{뒤집힌 \emph{출처}는 대상 라벨 없이 찾을 수 있다 --- 나머지 출처들의
합의와 반대면 그것이다. 일곱 대상 전부에서 애니를 정확히 지목한다.}{뒤집힌 것이
\emph{대상}이면 여섯이 다 같이 틀린 쪽을 보므로 합의로는 못 찾는다.}

\section{찾아도 고칠 것이 별로 없다}

\begin{table}[t]
\centering\tsize
\caption{균등 대비 짝지은 차이. 붓스트랩 200회.}
\begin{tabular}{lrl}
\toprule
& $\Delta$ & 95\% 구간 \\
\midrule
합의 뒤집기 & $+$0.0064 & $[-0.0085, +0.0213]$ \\
소프트 가중 & $+$0.0066 & $[-0.0070, +0.0238]$ \\
클립 가중 & $+$0.0073 & $[-0.0035, +0.0214]$ \\
\bottomrule
\end{tabular}
\end{table}

셋 다 보류다. 노트 72가 이미 이유를 냈다 --- 앙상블이 나쁜 출처를 여섯 중
하나로 희석하므로 고칠 손상이 이미 절반이다.

\begin{table}[t]
\centering\tsize
\caption{팝업 --- 뒤집힌 출처를 어떻게 할까.}
\begin{tabular}{lr}
\toprule
\textbf{그냥 둔다(균등 여섯)} & \textbf{$+$0.4594} \\
뺀다(다섯) & $+$0.4551 \\
뒤집는다 & $+$0.4348 \\
\bottomrule
\end{tabular}
\end{table}

\claimx{음의 상관 구성원을 그냥 두는 것이 뒤집거나 빼는 것보다 나을 수 있다 ---
팝업이 그렇다. 그 구성원의 오차가 나머지 다섯의 상관된 오차를 상쇄하기
때문이다.}{팝업 하나에서 본 것이고 $n{=}75$다. 아이돌에서는 반대로 뒤집는 쪽이
$+$0.064 낫다.}

\section{여섯 결정이 하나가 된다}

\begin{table}[t]
\centering\tsize
\caption{유보 평가 앙상블 $\rho$. 복제 150회.}
\begin{tabular}{lrrr}
\toprule
$k$ & 셀마다 부호(6번) & 합의$+$전역(1번) \\
\midrule
0 & $+$0.2748 & \textbf{$+$0.2800} \\
16 & $+$0.2945 & $+$0.2888 \\
24 & $+$0.3223 & $+$0.3206 \\
32 & \textbf{$+$0.3367} & $+$0.3327 \\
\bottomrule
\end{tabular}
\end{table}

$k{=}32$에서 $+$0.3327 대 $+$0.3367 --- 사실상 같다. \textbf{결정을 여섯 번에서
한 번으로 줄여도 같은 자리다.} 성적 이득은 아니지만 절차가 단순해지고, $k{=}0$
쪽은 합의가 $+$0.005 앞선다(공짜).

\section{게이트를 만들었다}

다른 여섯 대상에서는 정답을 안다 --- 그 대상들은 라벨이 있으므로 어느 쌍이
잘 되는지 이미 안다. 거기서 배워 새 대상에 옮기면 대상 라벨이 0건이다.
혼합 전문가(Jacobs 1991)의 게이팅 망과 같은 자리이고, 학습 표본이 셀 42개뿐이라
망 대신 능형 회귀를 쓴다. 대상 $t$의 가중치를 정할 때 $t$가 대상인 셀은 학습에서
전부 뺀다.

\begin{table}[t]
\centering\tsize
\caption{게이트. 대상 라벨 0건.}
\begin{tabular}{lrr}
\toprule
특징 집합 & 판정치 & 웹툰 \\
\midrule
\textbf{전체 8개} & \textbf{$+$0.2473} & \textbf{$+$0.2259} \\
대상 특징 제외 6개 & $+$0.2836 & $+$0.4806 \\
합의$+$잔차 & $+$0.2835 & $+$0.4782 \\
합의만 & $+$0.2838 & $+$0.4792 \\
\midrule
균등(기준) & $+$0.2744 & $+$0.4818 \\
\bottomrule
\end{tabular}
\end{table}

\textbf{전체 특징으로는 진다.} 그리고 웹툰 하나에서 $+$0.4818이 $+$0.2259로
무너진다.

\claimx{도메인 수준 특징(대상 $n$ · 대상 축 수)은 한 대상 안에서 상수이므로
게이트에 정보를 안 주고 자유도만 먹는다. 셀 42개로 보이는 학습 표본이 실제로는
\textbf{도메인 6개}다.}{능형 회귀 하나로 본 것이고 벌점을 조율하지 않았다.
더 센 정규화로는 다를 수 있다.}

특징을 빼면 게이트가 합의로 되돌아간다($+$0.2836 대 합의 클립 $+$0.2836).
\textbf{합의 가중치가 지금 데이터에서 배울 수 있는 것의 전부다.}

\begin{table}[t]
\centering\tsize
\caption{게이트가 본 것(표준화 계수, 전체 42셀 적합, $R^2{=}0.587$).}
\begin{tabular}{lr}
\toprule
공통 축 수 & $-$0.137 \\
출처 축 수 & $-$0.120 \\
log 대상 $n$ & $-$0.120 \\
대상 축 수 & $-$0.116 \\
합의 가중치 & $+$0.114 \\
정렬 잔차 & $-$0.024 \\
출처 자기 $\rho$ & $+$0.003 \\
\bottomrule
\end{tabular}
\end{table}

$R^2{=}0.587$이 특징 8개에 표본 42개다. 상위 넷 중 셋이 도메인 수준이라 사실상
``어느 도메인인가''를 외운 것이다.

\section{한계}

\begin{itemize}[leftmargin=1.1em,itemsep=0.15em]
\item 합의 정향도 게이트도 \textbf{채택하지 않았다}. 진단 도구로만 남긴다.
\item 합의는 뒤집힌 것이 \emph{대상}일 때 무력하다. 애니가 대상인 셀은 $-$0.277
      그대로다.
\item 게이트를 능형 회귀 하나로만 시도했다. 망도, 벌점 조율도, 비선형도 안 했다.
\item 팝업에서 ``그냥 두는 것이 낫다''를 $n{=}75$로 봤다. 다양성 효과인지
      우연인지 못 가른다.
\item 노트 72의 셀마다 부호와 여기의 합의$+$전역이 $k{=}32$에서 0.004 차이인데,
      복제 150회의 잡음이 그보다 크다. ``같다''고만 말할 수 있다.
\end{itemize}

\section{다음}

\begin{enumerate}[leftmargin=1.3em,itemsep=0.15em]
\item \textbf{여덟째 도메인.} 값이 여기서 분명해졌다 --- 셀 42$\to$56, 게이트
      유효 표본 6$\to$7, 합의 투표자 6$\to$7.
\item 쌍 수준 특징을 더 만든다 --- 도메인 수준 특징이 쓸모없다면 남는 길은
      쌍의 성질을 더 재는 것이다.
\item 대상이 뒤집힌 경우를 합의로 잡을 방법 --- 출처들의 \emph{쌍끼리} 관계를
      보면 될지 모른다.
\item 팝업의 다양성 효과를 다른 대상에서 확인한다.
\end{enumerate}

\vspace{0.6em}
\hrule height 0.3pt
\vspace{0.5em}
{\footnotesize
\textbf{참고문헌}\\[0.3em]
Jacobs, R. A. et al. (1991). Adaptive mixtures of local experts.
\emph{Neural Computation}, 3(1), 79--87.\\[0.15em]
Wolpert, D. H. (1992). Stacked generalization. \emph{Neural Networks}, 5(2), 241--259.\\[0.15em]
Breiman, L. (1996). Stacked regressions. \emph{Machine Learning}, 24(1), 49--64.\\[0.15em]
Kuncheva, L. I., Whitaker, C. J. (2003). Measures of diversity in classifier
ensembles. \emph{Machine Learning}, 51(2), 181--207.\\[0.15em]
Shazeer, N. et al. (2017). Outrageously large neural networks: the
sparsely-gated mixture-of-experts layer. \emph{ICLR}.
}

\end{document}
